% 数学ノート集・共通テンプレート v2（積分論2026年度第1部の版面を参考）
% A5判。章内で定義・定理・例・注意等の番号を共通化。
% コンパイル： platex -interaction=nonstopmode math_note_template_v2.tex (2回)
%              dvipdfmx math_note_template_v2.dvi
\documentclass[a5paper,10pt,onecolumn]{jsarticle}
\usepackage[dvipdfmx,top=22mm,bottom=20mm,inner=19mm,outer=18mm,headheight=13pt,headsep=6mm]{geometry}
\usepackage{amsmath,amssymb,amsthm,mathtools}
\usepackage{enumitem,fancyhdr,titlesec}
\usepackage[dvipdfmx,bookmarks=true,colorlinks=false]{hyperref}
\usepackage{pxjahyper}

% ----- 差し替える文書情報 -----
\newcommand{\NoteTitle}{解析接続と一価性}
\newcommand{\NoteSeries}{数学ノート集}

% ----- 文字・組版 -----
% 参考PDFの質素な明朝＋Computer Modernの印象を引き継ぎながら，
% A5判で窮屈にならないよう本文の行間を少し確保する．
\linespread{1.13}
\setlength{\parindent}{1zw}
\setlength{\parskip}{0pt}
\setlength{\abovedisplayskip}{6pt plus 2pt minus 2pt}
\setlength{\belowdisplayskip}{6pt plus 2pt minus 2pt}
\setlength{\abovedisplayshortskip}{3pt plus 2pt minus 1pt}
\setlength{\belowdisplayshortskip}{4pt plus 2pt minus 1pt}
\setlist[enumerate]{topsep=3pt,itemsep=1pt,parsep=0pt,leftmargin=2.4zw}
\titleformat{\section}{\normalfont\normalsize\bfseries}{\thesection}{0.8zw}{}
\titlespacing*{\section}{0pt}{2.0ex plus .4ex}{.9ex}
\titleformat{\subsection}{\normalfont\normalsize\bfseries}{\thesubsection}{0.7zw}{}
\titlespacing*{\subsection}{0pt}{1.5ex plus .3ex}{.6ex}
\raggedbottom

% ----- 柱とページ番号：ページ上部、罫線なし -----
\fancypagestyle{mathnote}{%
  \fancyhf{}%
  \fancyhead[L]{\footnotesize \NoteSeries\ （\NoteTitle）}%
  \fancyhead[R]{\footnotesize \thepage}%
  \renewcommand{\headrulewidth}{0pt}%
  \renewcommand{\footrulewidth}{0pt}%
}
\pagestyle{mathnote}

% ----- 節内通し番号の定義・定理群 -----
\newtheoremstyle{mathnotetheorem}
 {4.5pt}{4.5pt}{\normalfont}{}%
 {\bfseries}{}{0.65zw}%
 {\thmname{#1}\thmnumber{ #2}\thmnote{（#3）}}
\theoremstyle{mathnotetheorem}
\newtheorem{theorem}{定理}[section]
\newtheorem{definition}[theorem]{定義}
\newtheorem{lemma}[theorem]{補題}
\newtheorem{proposition}[theorem]{命題}
\newtheorem{corollary}[theorem]{系}
\newtheorem{example}[theorem]{例}
\newtheorem{remark}[theorem]{注意}
\renewcommand{\proofname}{証明}
\renewcommand{\qedsymbol}{$\square$}
\makeatletter
\renewenvironment{proof}[1][\proofname]{%
  \par\pushQED{\qed}\normalfont\topsep4\p@\@plus2\p@\relax
  \trivlist\item[\hskip\labelsep\bfseries #1]\ignorespaces
}{\popQED\endtrivlist\@endpefalse}
\makeatother

% ----- 使うノートで適宜拡張する記号 -----
\newcommand{\R}{\mathbb{R}}
\newcommand{\norm}[1]{\lVert #1\rVert}
\newcommand{\abs}[1]{\lvert #1\rvert}

\begin{document}
\section*{読了後にできること}
本稿を読むための前提は，正則関数の恒等定理，複素曲線積分，単連結性と道のホモトピーである．
正則関数の芽と解析接続を定義し，モノドロミー定理を本稿で証明する．
本稿の具体的な到達課題は，一価性が破れる具体的な解析接続を追跡する．

\section{正則関数の芽}

以下，$U\subset\mathbb C$ を領域とする．
開集合上の正則関数の恒等定理，Cauchyの積分定理，
および道とホモトピーの基本事項を前提とする．

\begin{definition}
$a\in U$ とする．
$a$ の開近傍 $V$ 上の正則関数 $f$ と，
$a$ の開近傍 $W$ 上の正則関数 $g$ が
$a$ のある開近傍で一致するとき
$(V,f)\sim_a(W,g)$ と定める．
この同値類を $a$ における\emph{正則関数の芽}といい，
$[f]_a$ と書く．
\end{definition}

\begin{lemma}\label{lem:germ}
$\sim_a$ は同値関係である．
また二つの芽 $[f]_a,[g]_a$ が等しいことと，
ある $r>0$ が存在して
$D(a,r)$ 上で $f=g$ となることは同値である
（円板は両方の定義域に含まれるように取る）．
\end{lemma}
\begin{proof}
反射律と対称律は定義から明らかである．
$(V,f)\sim_a(W,g)$ と
$(W,g)\sim_a(Z,h)$ なら，
$a$ の二つの開近傍の共通部分上で
$f=g$ および $g=h$ が成立するので
$(V,f)\sim_a(Z,h)$ である．
従って推移律も成立する．
最後の主張は，任意の開近傍が $a$ を中心とする
十分小さい開円板を含むことから従う．
\end{proof}

\begin{proposition}[芽の局所的一意性]\label{prop:germ-unique}
二つの正則関数 $f,g$ が連結な開集合 $V$ 上で定義され，
ある $a\in V$ について $[f]_a=[g]_a$ なら
$f=g$ が $V$ 上で成り立つ．
\end{proposition}
\begin{proof}
芽が一致するので $f=g$ は $a$ の近傍で成立する．
従って $f-g$ の零点集合は集積点を持ち，
正則関数の恒等定理から $f-g\equiv0$ となる．
\end{proof}

\section{局所的な解析接続}

\begin{definition}
領域 $V_1,V_2$ 上の正則関数 $f_1,f_2$ に対し，
$V_1\cap V_2$ のある連結成分 $W$ が空でなく，
$f_1|_W=f_2|_W$ が成立するとき，
$f_2$ は $W$ を通じた $f_1$ の\emph{解析接続}であるという．
\end{definition}

\begin{lemma}[重なりでの一意性]\label{lem:overlap}
$V_1\cap V_2$ の連結成分 $W$ において，
$f_1,f_2$ が一点の近傍で一致すれば
$W$ 全体で一致する．
\end{lemma}
\begin{proof}
$f_1-f_2$ は $W$ 上で正則であり，
零点集合は集積点を持つから恒等定理による．
\end{proof}

\begin{example}[幾何級数の接続]
$D(0,1)$ 上の正則関数
$f(z)=\sum_{n=0}^\infty z^n$ は
$1/(1-z)$ に等しい．
従って $\mathbb C\setminus\{1\}$ 上の
$F(z)=1/(1-z)$ は $f$ の解析接続を与える．
\end{example}
\begin{proof}
$|z|<1$ なら有限等比級数の和
$\sum_{n=0}^Nz^n=(1-z^{N+1})/(1-z)$ である．
$N\to\infty$ として $f(z)=1/(1-z)$ となる．
$F$ は $\mathbb C\setminus\{1\}$ 上で正則なので接続である．
\end{proof}

\section{曲線に沿う解析接続}

\begin{definition}
連続な曲線 $\gamma:[0,1]\to U$ と，
始点 $a=\gamma(0)$ の正則関数の芽 $g_0$ を取る．
$g_0$ が $\gamma$ に沿って\emph{解析接続可能}であるとは，
分割
$0=t_0<t_1<\cdots<t_N=1$ と
各 $\gamma([t_{j-1},t_j])$ を含む開円板
$D_j\subset U$ および正則関数 $f_j\in\mathcal O(D_j)$ が存在して，
\[
 [f_1]_a=g_0,\qquad
 [f_j]_{\gamma(t_j)}=[f_{j+1}]_{\gamma(t_j)}
 \quad(1\le j<N)
\]
が成り立つことをいう．
このとき終点の芽 $[f_N]_{\gamma(1)}$ を
接続によって得られた芽という．
\end{definition}

\begin{remark}
上の分割では $\gamma([t_{j-1},t_j])\subset D_j$ とするが，
そのような一つの円板が各短い部分曲線を含むように取れるのは，
連続曲線の一様連続性と $U$ の開性による．
既に与えられた解析接続の連鎖を細分しても
終点の芽は変わらない．
\end{remark}

\begin{lemma}[固定曲線に沿う接続の一意性]
\label{lem:path-unique}
同じ始点の芽を同じ曲線 $\gamma$ に沿って解析接続したとき，
終点に得られる芽は一意である．
\end{lemma}
\begin{proof}
二つの接続の連鎖を取り，両方の分割点を合わせて
共通細分を作る．始点の芽は一致する．
各小区間で対応する二つの正則関数は，
その小区間の始点の近傍で同じ芽を表すから，
共通部分に含まれる小さい円板上で一致する．
恒等定理によって小区間に沿う連続した重なりでも
同じ解析接続を与える．
これを有限個の小区間に対して帰納的に繰り返すと，
終点における芽が一致する．
\end{proof}

\begin{lemma}[経路の小さな変形に対する安定性]
\label{lem:path-stability}
始点と終点をそれぞれ固定した曲線 $\gamma$ に沿って，
芽 $g_0$ が解析接続可能であるとする．
このとき $\gamma$ の十分小さい一様近傍にある曲線
$\eta$ のうち始点・終点が同じものに沿っても
$g_0$ は解析接続可能であり，
終点の芽は $\gamma$ の場合と一致する．
\end{lemma}
\begin{proof}
$\gamma$ に沿う解析接続の有限連鎖
$(D_j,f_j,t_j)$ を取る．
必要なら分割を細かくし，
各閉部分曲線
$\gamma([t_{j-1},t_j])$ が $D_j$ のコンパクト部分集合に
含まれ，分割点 $\gamma(t_j)$ の周囲に
$f_j,f_{j+1}$ が一致する小円板があるようにする．
有限個のコンパクト集合と開円板の間には
正の余裕が取れる．
従って $\sup_t|\eta(t)-\gamma(t)|<\delta$ となる
十分小さい $\delta>0$ に対して，
$\eta([t_{j-1},t_j])\subset D_j$ が全ての $j$ で成り立つ．
また $\eta(t_j)$ は $f_j=f_{j+1}$ が成立する
重なりの小円板に入る．
従って同じ $f_j$ が $\eta$ に沿う解析接続の連鎖をなす．
始点と終点は固定しているため，
両端での芽は $\gamma$ の場合と一致する．
\end{proof}

\section{モノドロミー定理}

\begin{theorem}[ホモトピー不変性]\label{thm:homotopy}
$g_0$ を $a\in U$ における正則関数の芽とし，
$g_0$ は $a$ から始まる $U$ 内の
\emph{任意の連続曲線に沿って解析接続可能}であるとする．
始点 $a$ と終点 $b$ を固定した二つの曲線
$\gamma_0,\gamma_1$ が端点固定ホモトピーで結ばれるなら，
両曲線に沿って得られる終点の芽は一致する．
\end{theorem}
\begin{proof}
端点固定ホモトピーを
$H:[0,1]\times[0,1]\to U$ とし，
$\gamma_s(t)=H(s,t)$ と置く．
全ての曲線について解析接続が可能であるから，
各 $s$ に対して終点の芽 $G(s)$ が一意に定まる．
任意の $s_0$ について，
補題\ref{lem:path-stability}を $\gamma_{s_0}$ に適用する．
$H$ の連続性と定義域のコンパクト性から
$s\to s_0$ で
$\sup_t|\gamma_s(t)-\gamma_{s_0}(t)|\to0$
となる．
従って $s$ が $s_0$ に十分近ければ
$G(s)=G(s_0)$ である．
すなわち写像 $G:[0,1]\to
 \{\text{$b$ における正則関数の芽}\}$
は局所定数である．
区間 $[0,1]$ は連結なので $G$ は定数であり，
$G(0)=G(1)$ を得る．
\end{proof}

\begin{theorem}[モノドロミー定理]\label{thm:monodromy}
$U\subset\mathbb C$ を単連結な領域，
$a\in U$，$g_0$ を $a$ における正則関数の芽とする．
$g_0$ が $a$ から始まる任意の連続曲線に沿って
解析接続可能なら，
唯一の正則関数 $F:U\to\mathbb C$ が存在して
$[F]_a=g_0$ を満たす．
この $F$ の各点での芽は，
$g_0$ をその点へ至る任意の曲線に沿って
解析接続して得られる芽と一致する．
\end{theorem}
\begin{proof}
各 $z\in U$ に対して，
$a$ から $z$ に至る曲線 $\gamma_z$ を選ぶ．
$U$ は領域だから道連結であり，そのような曲線が存在する．
単連結性によって，同じ端点を持つ二つの曲線は
端点固定ホモトピーで結ばれる．
定理\ref{thm:homotopy}から，
$\gamma_z$ に沿う解析接続で得られる終点の芽は
経路の選択に依存しない．
その芽を $G_z$ と書き，
$F(z)$ を芽 $G_z$ の任意の代表関数の $z$ における値とする．
芽の定義からこの値はwell-definedである．

$z_0\in U$ を固定する．
$G_{z_0}$ は $z_0$ の小円板 $D\subset U$ 上の
正則関数 $h$ によって表される．
任意の $z\in D$ について，
$a$ から $z_0$ への経路に
$D$ 内の線分 $[z_0,z]$ を連結すると，
その経路に沿う解析接続の終点の芽は $[h]_z$ になる．
経路独立性から $G_z=[h]_z$ であり，
$F(z)=h(z)$ が $D$ 全体で成立する．
従って $F$ は $z_0$ の近傍で正則であり，
$z_0$ は任意なので $F$ は $U$ 上正則である．
$a$ における芽が $g_0$ であることは定義から従う．
最後に別の正則関数 $F'$ が同じ芽を持つなら，
恒等定理から $F'=F$ であり一意である．
\end{proof}

\begin{remark}
単連結性だけでは結論を導けない．
全ての曲線に沿って解析接続できるという仮定が本質的である．
例えば $U=\mathbb C$ として原点以外で
$1/z$ を表す芽を取っても，
原点を終点とする曲線に沿って正則な接続はできない．
\end{remark}

\section{複素対数と解析接続}

\begin{lemma}[局所対数]\label{lem:local-log}
$a\ne0$ に対し，$a$ の十分小さい円板 $D$ 上に
正則関数 $\ell$ が存在して
$e^{\ell(z)}=z$ を満たす．
\end{lemma}
\begin{proof}
$a$ の対数値 $c$ を一つ選び $e^c=a$ とする．
$|z-a|<|a|$ で
\[
 \ell(z)=c+\sum_{n=1}^\infty
 \frac{(-1)^{n+1}}n\left(\frac{z-a}{a}\right)^n
\]
と置く．級数は開円板で局所一様収束し，
項別微分によって
$\ell'(z)=1/z$ となる．
$h(z)=e^{\ell(z)}/z$ と置くと
$h'/h=\ell'-1/z=0$ であり，
$h(a)=1$ なので連結性から $h\equiv1$ となる．
従って $e^{\ell(z)}=z$ である．
\end{proof}

\begin{proposition}[対数の解析接続と巻き数]\label{prop:log-winding}
$\gamma:[0,1]\to\mathbb C^\times$ を
始点 $a$ と終点 $b$ を持つ区分的 $C^1$ 曲線とする．
$a$ の対数の芽 $\ell_0$ を $\gamma$ に沿って接続すると，
終点で得られる芽 $\ell_1$ は
\[
 \ell_1(b)-\ell_0(a)=\int_\gamma\frac{dz}{z}
\]
を満たす．
特に $\gamma$ が閉曲線なら，
接続後の対数値は元の対数値に
$2\pi i\operatorname{Ind}(\gamma,0)$ を加えたものとなる．
\end{proposition}
\begin{proof}
局所対数は導関数が $1/z$ である．
$\gamma$ の像はコンパクトなので，
有限個の局所対数の円板で覆い，
隣接する対数の値が分割点で一致するように
$2\pi i$ の整数倍を調整できる．
各小区間で複素積分の基本定理を適用すると，
局所対数の終点と始点の値の差は
$\int dz/z$ に等しい．
これらを全区間について足し合わせると
中間点の値が打ち消し合い，
最初の表示式を得る．
閉曲線なら
$\int_\gamma dz/z=2\pi i\operatorname{Ind}(\gamma,0)$
は巻き数の定義とその整数性による．
\end{proof}

\begin{example}[原点を一周する対数]
$\gamma(t)=e^{2\pi it}$（$0\le t\le1$）とする．
$1$ における対数の芽を $\gamma$ に沿って接続すると，
終点での対数値は始点の対数値より $2\pi i$ 大きい．
従って $\mathbb C^\times$ 上に
一価正則な複素対数は存在しない．
\end{example}
\begin{proof}
\[
 \int_\gamma\frac{dz}{z}
 =\int_0^1\frac{2\pi ie^{2\pi it}}{e^{2\pi it}}\,dt
 =2\pi i
\]
なので命題\ref{prop:log-winding}による．
もし $F$ が $\mathbb C^\times$ 上の正則対数なら
$F'(z)=1/z$ だから，閉曲線上のその積分は
原始関数の存在によって零となるはずであり矛盾する．
\end{proof}

\begin{theorem}[単連結領域上の正則対数]\label{thm:global-log}
$U$ を単連結な領域，
$f:U\to\mathbb C^\times$ を零を持たない正則関数とする．
このとき正則関数 $g:U\to\mathbb C$ が存在して
\[
 e^{g(z)}=f(z)\qquad(z\in U)
\]
を満たす．$g$ は $2\pi i$ の整数倍の定数を除いて一意である．
\end{theorem}
\begin{proof}
$f'/f$ は $U$ 上で正則である．
単連結領域上のCauchy積分定理から，
正則関数 $h$ で $h'=f'/f$ となるものが存在する．
ある $a\in U$ を固定し，$e^c=f(a)$ となる $c$ を選ぶ．
\[
 g(z)=h(z)-h(a)+c
\]
と置けば $g'=f'/f$ であり，
$\frac d{dz}(e^g/f)=0$ となる．
$e^{g(a)}/f(a)=1$ なので
$e^{g(z)}=f(z)$ が全域で成り立つ．
二つの正則対数 $g_1,g_2$ は
$e^{g_1-g_2}=1$ を満たす．
従って $g_1-g_2$ は離散集合 $2\pi i\mathbb Z$ に値を取る
連続関数であり，$U$ が連結だから定数である．
\end{proof}

\section{平方根と分岐}

\begin{proposition}[平方根の解析接続]\label{prop:sqrt}
$\mathbb C^\times$ 上の局所正則な平方根の芽を
原点の周りに一周解析接続すると，
その値は元の値の符号を反転する．
従って $\mathbb C^\times$ 上には
$g(z)^2=z$ を満たす一価正則関数 $g$ は存在しない．
\end{proposition}
\begin{proof}
局所対数 $\ell$ によって
$g(z)=e^{\ell(z)/2}$ と書ける．
原点を一周すると
命題\ref{prop:log-winding}から
$\ell$ は $2\pi i$ 増える．
従って $g$ は
$e^{(\ell+2\pi i)/2}=-e^{\ell/2}$ へ変わる．
もし一価正則な平方根があれば，
解析接続はその関数自身によって与えられるため
一周後の芽も同じとなるが，符号の反転と矛盾する．
\end{proof}

\begin{definition}
局所的な解析接続の全ての芽を集め，
解析接続可能な関係で結んだ空間を考えると，
元の領域の上に射影する空間が得られる．
対数や平方根の場合，この空間は
多価関数を一価な関数として表現する
\emph{Riemann面}の基本例となる．
\end{definition}

\begin{remark}
モノドロミー定理は解析接続に関する
大域的一価性を，領域の単連結性と
経路に沿う接続可能性から導く定理である．
多価性は単に関数の表記上の問題ではなく，
閉曲線に沿う解析接続が芽へ与える作用に現れる．
\end{remark}

\section{到達課題の完全解答}
\noindent\textbf{対数関数の一周による値の変化．}
\begin{proposition}
$\mathbb C^\times$ における局所正則対数の枝を，
原点を正向きに一周する閉曲線に沿って解析接続すると，
終点での値は始点の値より $2\pi i$ 増加する．
\end{proposition}
\begin{proof}
$\gamma(t)=e^{it}$（$0\le t\le2\pi$）とする．
$\gamma$ に沿った正則対数の局所値 $L(t)$ は
$e^{L(t)}=\gamma(t)$ を満たす．
微分すると
$L'(t)=\gamma'(t)/\gamma(t)=i$ なので
$L(2\pi)-L(0)=\int_0^{2\pi}i\,dt=2\pi i$．
従って原点を囲む環状領域では
局所対数を一価な正則関数へ貼り合わせることはできない．
\end{proof}
\begin{remark}
この計算は，局所的な正則性と
大域的な一価性が異なる条件であることを表す．
\end{remark}
\noindent\textbf{到達確認．}
同じ曲線を逆向きに一周した場合の変化量を求めよ．
\par\noindent\textbf{解答．}
逆向きでは $\int dz/z=-2\pi i$ なので
対数の解析接続による変化量は $-2\pi i$ である．
\end{document}